% ---
% title: Bond defaults: independent or not?
% tags: finance, independence, inclusion-exclusion, complement, conditional-probability
% ---

\begin{problem}
	\begin{context}
		A bond \emph{defaults} when the issuer fails to pay what it promised. A
		\emph{first-to-default swap} on a basket of bonds is an insurance-like
		contract that pays as soon as at least one bond of the basket defaults;
		a \emph{second-to-default swap} pays only when a second bond defaults as
		well. Their prices depend on how the defaults are related: issuers that
		depend on the same economy tend to get into trouble together. This is
		called \emph{default correlation}.
	\end{context}
	A basket contains two bonds. Let $D_1$ and $D_2$ be the events that bond 1
	and bond 2 default within a year; $P(D_1)=P(D_2)=0.05$. Find the
	probability that at least one bond defaults (then a first-to-default swap
	pays) and the probability that both bonds default (then a
	second-to-default swap pays),
	(a) if the two defaults are independent;
	(b) if, because both issuers depend on the same economy, the probability
	that both default is $P(D_1D_2)=0.02$.
	(c) In case (b), find $P(D_2\mid D_1)$ and compare it with $P(D_2)$. Are
	the defaults independent?
	(d) Back to independence: a basket contains 10 bonds that default
	independently of each other, each with probability $0.05$. What is the
	probability that at least one of them defaults within a year?
\end{problem}

\begin{slideproblem}
	Two bonds, $D_i$ = ``bond $i$ defaults within a year'',
	$P(D_1)=P(D_2)=0.05$. Find $P(\text{at least one defaults})$ (a
	first-to-default swap pays) and $P(\text{both default})$ (a
	second-to-default swap pays),
	(a) if the defaults are independent;
	(b) if $P(D_1D_2)=0.02$ (same economy).
	(c) In case (b): $P(D_2\mid D_1)=\,$? Compare with $P(D_2)$.
	(d) 10 bonds default independently, each with probability $0.05$.
	$P(\text{at least one default})=\,$?
\end{slideproblem}

\begin{slides}
	\small
	\exstatement
	\pause
	\sol ``At least one defaults'' $=D_1\cup D_2$; \ ``both default''
	$=D_1D_2$.
	\pause\\[1mm]
	(a) Independent: the probability of both is the product.
	\[ P(D_1D_2)=P(D_1)P(D_2)=0.05\cdot0.05=\alert{0.0025}. \]
	\pause
	\vspace{-2mm}
	\[ P(D_1\cup D_2)=P(D_1)+P(D_2)-P(D_1D_2)=0.05+0.05-0.0025
		=\alert{0.0975}. \]
	\pause
	Check: no default $=D_1^cD_2^c$, also independent:
	$1-0.95^2=0.0975$.
\end{slides}

\begin{slides}
	\small
	(b) Same economy: $P(D_1D_2)=\alert{0.02}$ is given.
	\pause\\[1mm]
	This is 8 times as large as $0.0025$: no product rule here.
	\pause
	\[ P(D_1\cup D_2)=0.05+0.05-0.02=\alert{0.08}. \]
	\pause
	(c) Definition of conditional probability:
	\[ P(D_2\mid D_1)=\frac{P(D_1D_2)}{P(D_1)}=\frac{0.02}{0.05}=\alert{0.4}. \]
	\pause
	$0.4\ne0.05=P(D_2)$: the defaults are \alert{not independent}.
	\pause\\[1mm]
	Once bond 1 has defaulted, bond 2 defaults with probability $0.4$
	instead of $0.05$: \emph{positive} dependence.
\end{slides}

\begin{slides}
	\small
	(d) Complement: ``at least one default'' is the opposite of ``all 10
	bonds survive''.
	\pause\\[1mm]
	Each bond survives with probability $0.95$; independent, so multiply:
	\[ P(\text{no default})=0.95^{10}\approx0.599. \]
	\pause
	\vspace{-2mm}
	\[ P(\text{at least one default})=1-0.95^{10}\approx\alert{0.401}. \]
	\pause
	Each bond is fairly safe, but with 10 of them a first-to-default swap
	pays with probability about 40\%.
	\pause
	\begin{center}
		\begin{tabular}{lcc}
			two bonds    & at least one & both     \\ \hline
			independent  & $0.0975$     & $0.0025$ \\
			same economy & $0.08$       & $0.02$
		\end{tabular}
	\end{center}
	\alert{Lesson:} positive dependence makes ``at least one'' slightly less
	likely, but ``both'' much more likely.
\end{slides}

\begin{solution}
	\idea Write the two swap events as $D_1\cup D_2$ and $D_1D_2$. Under
	independence the probability of the intersection is the product; then the
	union follows from inclusion--exclusion. Without independence the
	intersection must be given. For many bonds, use the complement.

	\step{Events.} The first-to-default swap pays if at least one bond
	defaults: this is the union $D_1\cup D_2$. The second-to-default swap pays
	if both bonds default: this is the intersection $D_1D_2$. For the union we
	always have the inclusion--exclusion formula
	\[ P(D_1\cup D_2)=P(D_1)+P(D_2)-P(D_1D_2). \]
	So in every case the key number is $P(D_1D_2)$.

	\step{(a)} The defaults are independent. By definition this means that the
	probability of both is the product of the probabilities:
	\[ P(D_1D_2)=P(D_1)P(D_2)=0.05\cdot0.05=0.0025. \]
	Then inclusion--exclusion gives
	\[ P(D_1\cup D_2)=0.05+0.05-0.0025=0.0975. \]
	We subtract $0.0025$ because the outcomes where both bonds default are
	counted once in $P(D_1)$ and once more in $P(D_2)$.

	Check with the complement. ``At least one defaults'' is the complement of
	``no bond defaults'', which is $D_1^cD_2^c$ (De Morgan). If $D_1$ and $D_2$
	are independent, so are $D_1^c$ and $D_2^c$, and each has probability
	$1-0.05=0.95$. So
	\[ P(D_1\cup D_2)=1-P(D_1^cD_2^c)=1-0.95\cdot0.95=1-0.9025=0.0975. \]

	\step{(b)} Now the two issuers depend on the same economy, and
	$P(D_1D_2)=0.02$ is given. We cannot multiply here: multiplication is
	allowed only for independent events. This number is $0.02/0.0025=8$ times
	the value under independence. The probability that at least one bond
	defaults is again inclusion--exclusion:
	\[ P(D_1\cup D_2)=0.05+0.05-0.02=0.08. \]

	\step{(c)} By the definition of conditional probability,
	\[ P(D_2\mid D_1)=\frac{P(D_1D_2)}{P(D_1)}=\frac{0.02}{0.05}=0.4. \]
	Two events are independent exactly when knowing that one happened does not
	change the probability of the other, that is, $P(D_2\mid D_1)=P(D_2)$.
	Here $0.4\ne0.05$, so the defaults are not independent. The dependence is
	positive: once bond 1 has defaulted (probably because the economy is in
	trouble), bond 2 defaults with probability $0.4$, eight times its
	unconditional probability $0.05$. Equivalently,
	$P(D_1D_2)=0.02\ne0.0025=P(D_1)P(D_2)$.

	\step{(d)} With 10 bonds, ``at least one default'' would need
	inclusion--exclusion with many terms. The complement is much simpler: ``at
	least one bond defaults'' is the complement of ``no bond defaults'', that
	is, of ``all 10 bonds survive''. Each bond survives with probability
	$1-0.05=0.95$, and the survivals are independent (as the defaults are), so
	the probability that all 10 survive is the product of 10 equal factors:
	\[ P(\text{no default})=0.95^{10}\approx0.5987. \]
	Therefore
	\[ P(\text{at least one default})=1-0.95^{10}\approx1-0.5987=0.4013. \]
	Although each bond is fairly safe, a first-to-default swap on 10 such
	bonds pays with probability about 40\%.
	\begin{answer}
		(a) $P(D_1\cup D_2)=0.0975$, $P(D_1D_2)=0.0025$ \quad
		(b) $P(D_1\cup D_2)=0.08$, $P(D_1D_2)=0.02$ \quad
		(c) $P(D_2\mid D_1)=0.4\ne0.05=P(D_2)$: not independent \quad
		(d) $1-0.95^{10}\approx0.401$
	\end{answer}
	\begin{remark}
		Compare (a) and (b). Positive dependence makes ``at least one default''
		slightly \emph{less} likely ($0.08$ instead of $0.0975$), but ``both
		default'' much \emph{more} likely ($0.02$ instead of $0.0025$). The
		defaults tend to come together, so they overlap more and cover less of
		$\Om$. This is why the seller of a second-to-default swap cares a lot
		about default correlation.
		\textbf{Common mistake:} computing $P(D_1D_2)=0.05\cdot0.05$ in (b)
		as well. The product rule holds only for independent events.
	\end{remark}
\end{solution}
